% Platonic-ideal introduction (Chapter 1) opening.
% Author: Raeez Lorgat. Chriss--Ginzburg register. No AI attribution.

\providecommand{\ClaimStatusTheorem}{\textsc{[Theorem]}}
\providecommand{\ClaimStatusConjectured}{\textsc{[Conjectured]}}
\providecommand{\ClaimStatusDefinition}{\textsc{[Definition]}}
\providecommand{\kBKM}{\kappa_{\mathrm{BKM}}}
\providecommand{\kch}{\kappa_{\mathrm{ch}}}
\providecommand{\kcat}{\kappa_{\mathrm{cat}}}
\providecommand{\kfib}{\kappa_{\mathrm{fiber}}}
\providecommand{\CYcat}{\mathrm{CY}\text{-cat}}
\providecommand{\ChirAlg}{\mathrm{ChirAlg}}
\providecommand{\HolFA}{\mathrm{HolFA}}
\providecommand{\hCS}{\mathrm{hCS}}
\providecommand{\Sp}{\mathrm{Sp}}
\providecommand{\cF}{\mathcal{F}}
\providecommand{\cE}{\mathcal{E}}
\providecommand{\cO}{\mathcal{O}}
\providecommand{\cM}{\mathcal{M}}
\providecommand{\fg}{\mathfrak{g}}
\providecommand{\CoHA}{\mathrm{CoHA}}
\providecommand{\Meff}{\mathcal{M}^{+}_{\mathrm{eff}}}
\providecommand{\ZZ}{\mathbb{Z}}
\providecommand{\CC}{\mathbb{C}}
\providecommand{\bH}{\mathbf{H}}

% Lossless insertion directive: this file is a new Section~1 of
% chapters/theory/introduction.tex. The existing introduction chapter
% title remains. Existing sections shift from $n$ to $n+1$ in numbering;
% no section is removed, rewritten, or shortened.

\section{The positive-geometry grammar: the object of study}
\label{wn:sec:platonic-intro-object-wn}

Let $\cA$ be a Calabi--Yau category of complex dimension $d$ in the
Kontsevich--Soibelman sense: a pretriangulated differential graded
category equipped with a non-degenerate cyclic pairing of degree $-d$
respecting the $A_\infty$-structure up to coherent homotopy. Such
a category sits on a Calabi--Yau variety $X$ in one of two
complementary senses. Either $X$ carries $\cA$ as the derived category
of coherent sheaves $D^b \mathrm{Coh}(X)$ equipped with its Serre-duality
cyclic structure; or $X$ carries $\cA$ as a category of
matrix factorisations of a Landau--Ginzburg potential $W$ on a
Calabi--Yau quiver with potential $(Q, W)$. The Van~den Bergh
non-commutative crepant resolution bridges the two: any smooth
affine Calabi--Yau threefold admits a presentation as the bounded
derived category of a Calabi--Yau algebra $A_Q = \mathbb{C}[Q] /
\partial W$; any compact Calabi--Yau threefold admits a cover by
such charts glued by Morita tilting cocycles.

Write $\Omega_X \in H^{d, 0}(X)$ for the holomorphic volume form. This
volume form is the load-bearing datum: all structure derives from it.

\subsection{The functor}
\label{wn:sec:platonic-intro-functor-wn}

The central object of the present volume is the Calabi--Yau-to-chiral
functor
\[
\Phi_d: \CYcat_d \longrightarrow \ChirAlg.
\]
The preface announces its canonical factorisation. The introduction
states its components.

\begin{theorem}[Stage 1: the holomorphic factorisation algebra]
\label{wn:thm:platonic-intro-stage-1}
\ClaimStatusTheorem\
There is a canonical functor
\[
\Phi^{\mathrm{FA}}_d: \CYcat_d \longrightarrow E_d\text{-}\HolFA(X)
\]
sending a Calabi--Yau category $\cA$ over $X$ to the
Costello--Gwilliam holomorphic factorisation algebra of observables of
the six-dimensional (resp.\ $2d$-dimensional) holomorphic Chern--Simons
theory on $X$ with gauge data read from $\cA$. The assignment is
determined up to contractible choice by: the Kontsevich 1999 /
Tamarkin 2003 $E_d$-formality of the Dolbeault little-$d$-disks
operad; the Costello--Gwilliam 2017 factorisation axiom; the
Costello--Li 2016 holomorphic-twist identification at $d = 3$; the
Francis--Gaitsgory 2012 chiral Koszul duality for $E_d$-algebras;
and the Gwilliam--Williams 2021 $E_d^{\mathrm{hol}}$-to-$E_d$ comparison.
\end{theorem}

\begin{theorem}[Stage 2: specialisation to a curve]
\label{wn:thm:platonic-intro-stage-2}
\ClaimStatusTheorem\
For any closed $(d{-}1)$-dimensional subvariety $\Sigma_{d-1} \subset X$
and reference curve $C \subset X$,
factorisation homology fibrewise along the projection
$X \to C$ yields a functor
\[
\Sp_{\Sigma_{d-1}, C}: E_d\text{-}\HolFA(X) \longrightarrow
E_1\text{-}\ChirAlg(C).
\]
The composite $\Phi_d = \Sp_{\Sigma_{d-1}, C} \circ \Phi^{\mathrm{FA}}_d$
depends on the pair $(\Sigma_{d-1}, C)$. Different choices produce
different chiral-algebra shadows of the same Stage-1 output.
\end{theorem}

The six routes to $G(K3 \times E)$ catalogued in prior work
(Mukai-lattice, Igusa $\Phi_{10}$ via Gritsenko $\Delta_5$, BKM
Borcherds weight, K3 fibre rank, Niemeier lattice, Humbert divisor,
CHL twisted orbifold) are \emph{not} six applications of $\Phi_3$
to one object: they are six different kinds of measurement.
Four of them (Mukai, Hodge, fibre, Igusa-weight) are Stage-1
invariants of $\cF_\cA$ or of the underlying Calabi--Yau datum; the
others are $(\Sigma_2, C)$-specialisations of a single canonical
$\cF_\cA \in E_3\text{-}\HolFA(K3 \times E)$ at different cycle
choices. Reading the six faces through this grammar dissolves
their apparent conflict.

\subsection{The quantum group}
\label{wn:sec:platonic-intro-quantum-group-wn}

Inside $\cF_\cA$ lives the positive effective moduli stack
$\Meff(X)$ carrying the Kontsevich--Soibelman critical cohomological
Hall algebra structure. Its equivariant vanishing-cycle cohomology
is the positive half of the quantum group:
\[
Y^+(X) \;:=\; H^{\bullet}_{\mathrm{eq}}\!\bigl(\Meff(X),\,\phi_W\bigr).
\]
The Drinfeld double
\[
G(X) \;:=\; D(Y^+(X)) \;=\; Y^+(X) \bowtie Y^0(X) \bowtie Y^-(X)
\]
is the quantum group attached to $\cA$. The Hopf pairing comes from
Serre duality on the compact locus, supplemented by
Davison--Meinhardt PBW integration; the spectral structure comes from
Maulik--Okounkov stable envelopes transported by positive-half
gluing cocycles.

The equivariance label stratifies by geometric type: full toric $T^d$
on $\CC^3$ and toric three-folds, reduced virtual class with residual
$\CC^\times$ on compact non-toric Calabi--Yau threefolds, discrete
orbifold inertia on quotient geometries, period-domain equivariance
on lattice-polarised families. The algebra $Y^+(X)$ is the same
construction read through four local languages.

\subsection{The chiral shadow on a curve}
\label{wn:sec:platonic-intro-chiral-shadow-wn}

Specialising Stage 1 at $(\Sigma_2, C) = (K3, E)$ sends
$\cF_\cA \in E_3\text{-}\HolFA(K3 \times E)$ to an $E_1$-chiral
algebra on the elliptic curve $E$. Maulik--Okounkov stable-envelope
transport identifies this chiral algebra with the chiral bialgebra
$\bH_{\Delta_5}$ supported on the Humbert divisor $H_1$ inside the
Siegel paramodular space $\cH_{3,2}$: its character is $\Delta_5^{-2}$,
its root-multiplicity table is the Fourier coefficients of the
Gritsenko weak Jacobi form $\phi_{0,1}$, and its Weyl--Kac--Borcherds
denominator identity is the singular theta lift of $\phi_{0,1}$ in
the sense of Borcherds 1998.

The generalised Kac--Moody superalgebra $\fg_{\Delta_5}$ of
Gritsenko--Nikulin 1998 is the pull-back along the automorphic
correction; a programme-level conjecture identifies it with the
bracket completion of the BPS Lie algebra extracted from
$\Meff(K3 \times E)$ by Davison--Meinhardt integrality. The
graded dimensions coincide unconditionally; the bracket-level
identification is the canonical frontier, outlined in the
$K3 \times E$ Calabi--Yau threefold programme chapter
(\texttt{k3e\_cy3\_programme}).

\subsection{The eight-form class}
\label{wn:sec:platonic-intro-eight-form-wn}

Gritsenko and Cl\'ery 2008 classified the diagonal-divisor
paramodular cusp forms of Hecke type. There are exactly eight.
In weight-ordered list:
\[
\Delta_5, \quad \Delta_2^{(2)}, \quad \Delta_1^{(3)}, \quad
\Delta_1^{(4)}, \quad \Delta_{1/2}^{(5)}, \quad \Delta_1^{(6)},
\quad \Delta_{1/4}^{(7)}, \quad \Delta_0^{(8)}.
\]
The universal weight identity $\kBKM(\Phi_N) = c_N(0)/2$ holds
uniformly across all eight forms with cover selection
$\{\Sp_4, \Mp_4, \widetilde{\Mp}_4\}$ determined by weight
integrality. The Chaudhuri--Hockney--Lykken subset
$N \in \{1, 2, 3, 4, 6\}$ --- precisely those $N$ with
$\varphi(N) \mid 2$ --- realises each $\Phi_N$ as the denominator
of a generalised Kac--Moody superalgebra on the twisted-reduced
Donaldson--Thomas moduli of the Calabi--Yau orbifold
$(K3 \times E)/(\ZZ/N\ZZ)$; the three remaining $N \in \{5, 7, 8\}$
fall outside the CHL slice and find their geometric hosts in
the Borcea--Voisin threefold $(S_5 \times E_5)/(\iota_S \times \iota_E)$
at $N = 5$ with metaplectic fourth-root character, an
order-$4$ Cheeger--Simons gerbe on the singular
$(K3 \times E)/\ZZ_7$ at $N = 7$, and the Mongardi--Tari--Wandel
Kummer-$3$ hyperk\"ahler fourfold at $N = 8$.

Each form supplies a four-corner data: a Calabi--Yau geometric host,
a cohomological Hall algebra extracted from its positive effective
moduli, a BPS Lie algebra extracted by Davison--Meinhardt
integrality, and an automorphic correction realising the
generalised Kac--Moody as a Borcherds singular theta lift.
For $N = 1$ these four corners commute unconditionally (with
the bracket-level matching conditional); for $N \in \{2, 3, 4, 6\}$
they commute on the Bryan--Oberdieck primitive-class sublocus,
with the imprimitive completion conditional on Oberdieck--Shen;
for $N \in \{5, 7, 8\}$ they commute conditionally on the three
non-CHL host constructions listed above.

\subsection{Structure of the volume}
\label{wn:sec:platonic-intro-structure-wn}

Part~I (Foundations: \texttt{introduction} through
\texttt{hochschild\_calculus}) assembles the
$A_\infty / L_\infty$ algebraic apparatus, the Calabi--Yau category
definitions, and the Hochschild calculus.

Part~II (The functor $\Phi$: \texttt{cy\_to\_chiral} through
\texttt{drinfeld\_center}) constructs $\Phi^{\mathrm{FA}}_d$
and $\Sp_{\Sigma_{d-1}, C}$, proves the two-stage factorisation,
and develops the six-dimensional holomorphic Chern--Simons
machinery.

Part~III ($E_n$ hierarchy and dimension stratification:
\texttt{e1\_chiral} through \texttt{en\_factorization})
tracks how $\Phi_d$ lands in $E_1, E_2, E_d$-chiral algebras for
$d = 1, 2, \geq 3$ respectively.

Part~IV (Quantum groups and Drinfeld centre:
\texttt{quantum\_groups\_foundations} through
\texttt{drinfeld\_center}) constructs $G(X) = D(Y^+(X))$ and
catalogues its $R$-matrix, braided tensor category, and
$(\infty, 1)$-centre structures.

Part~V (The Calabi--Yau landscape:
\texttt{derived\_categories\_cy} through
\texttt{k3e\_cy3\_programme}) realises the functor on K3, on
$K3 \times E$, on the conifold, on local $\mathbb{P}^2$, on elliptic
Calabi--Yau threefolds, and on the eight-form class.

Part~VI (Seven faces of $r_{\mathrm{CY}}$:
\texttt{cy\_d\_kappa\_stratification} through
\texttt{quantum\_group\_reps}) catalogues the stratification of
$\Phi$-output by Calabi--Yau dimension and organises the six (or
seven) routes to $G(K3 \times E)$ into two-stage specialisations.

Part~VII (Frontiers: \texttt{modular\_trace} through the end)
houses the conjectural extensions, the non-CHL hosts, the
higher-dimensional cousins at $d = 5$, and the catalogue of open
problems.

The reader who wishes to see the full positive-geometry grammar
concretely should turn to the $K3 \times E$ Calabi--Yau threefold
programme (Part~V, \texttt{k3e\_cy3\_programme}) first, then
return to Part~II for the foundational theorems.
